\section{Fragmento de código de ejemplo de DSPy}

\subsection{Pipeline declarativo}

\begin{lstlisting}[language=Python]
from dspy import Signature, Module, Optimizer, Evidence

class RetrieveModule(Module):
    signature = Signature("question", "context")
    def run(self, question):
        return search_knn(question, top_k=5)

class ReasonModule(Module):
    signature = Signature("context", "answer")
    def run(self, context):
        return llm_chat(f"Use the context to answer: {context}")

class VerifyModule(Module):
    signature = Signature("answer", "rating")
    def run(self, answer):
        score = verificator(answer)
        return {"score": score, "flags": []}

pipeline = [
    RetrieveModule(),
    ReasonModule(),
    VerifyModule()
]

optimizer = Optimizer(pipeline)
optimizer.search(metric="score", budget=100)

evidence = Evidence(pipeline, optimizer.result)
evidence.save("./evidence/latest.json")

# Additional helper routine used in experiments
def plan_and_execute(goal):
    plan = optimizer.plan(goal)
    for step in plan.steps:
        evidence.log(f"Executing {step.name}")
        result = step.module.run(step.input)
        evidence.attach(step.name, result)
        if result.get("score", 100) < 70:
            evidence.mark_failure(step.name)
            break
    return evidence.finalize()

# More DSL helpers for debugging
def replay_evidence(snapshot_path):
    snapshot = Evidence.load(snapshot_path)
    for entry in snapshot.entries:
        print(f"> {entry.prompt[:80]}... → {entry.metric}")
        if entry.metric < 0.6:
            print("   Needs review: ", entry.flags)
\end{lstlisting}

\subsection{Resumen de la sección}

Este fragmento de código muestra la estructura declarativa de DSPy combinada con optimizers, y ayuda al equipo a entender la secuencia real de llamadas detrás del compilador.

